\documentclass[11pt]{article}
\usepackage[T1]{fontenc}
\usepackage{amsmath,amssymb,amsthm,geometry}
\geometry{margin=.9in}
\setlength{\emergencystretch}{2em}
\allowdisplaybreaks[2]
\newtheorem{theorem}{Theorem}
\newtheorem{lemma}{Lemma}
\newcommand{\E}{\mathbb E}
\newcommand{\dd}{\,d}
\title{The residual-budget barrier: two signed exchange channels}
\author{E5 continuation, analytic partial result}
\date{Version 1}
\begin{document}\maketitle

\paragraph{Disposition.}
The required barrier for the actual critical trajectory remains open.
We prove that, up to RQT's reached stopping event, the accumulated
exchange omitted by a four-regression reduction is $o(1)$.
Consequently a fixed-size residual exit would require a signed payment
from two explicit cross-output channels. We neither prove an exit nor
choose which channel dominates on the initialized path.

\section{Actual initialization and the exact face}
Use the actual aligned balanced isotropic flow with
$S_0=cq^{1/\omega}$, fixed $c>0$, $\rho\in[13/20,7/10]$,
$\lambda=\rho^2$, and $\omega=(1-\lambda)/2$.
The two populations are $P_1=N(e_1,q^2I)$ and
$P_2=N(\rho e_2,q^2I)$ with mixture weights $1/2$.
Write $f(D)=\int W(U^TD)_+\dd l$, $|U|=|W|$,
$S=\int|U|^2\dd l$, $r=D-f(D)$, and
$\mathcal L=\frac14\sum_p\E_{P_p}|r|^2$.
The parameter gradient uses both physical vectors $(U,W)$ and the
fixed probability label measure $\dd l$.
Put $Q_1=1+q^2$, $Q=\lambda+q^2$, and
\[
 m=\frac{\E_{P_1}D_1f_1}{Q_1},\qquad
 z=\frac{\E_{P_1}D_1f_2}{Q_1},\qquad
 h=\frac{\E_{P_2}D_2f_1}{Q},\qquad
 n=\frac{\E_{P_2}D_2f_2}{Q}.
\]
Here $z,h$ are signed cross regressions, and $n$ is the actual weak
response. None of these is identified with a cohort's radial mass.

RQT supplies an initialized stopping time. Choose its small hit
$\tau_0$, with $n(\tau_0)=\eta_0>0$, so that, for
$r_*=\lambda/800$,
\[
 \mathcal R(\tau_0)\le r_*^2/16,\qquad
 \mathcal R=4\mathcal L-Q(1-n)^2.
\]
Let $\tau$ be the first subsequent hit of either $n=1/2$ or
$\mathcal R=r_*^2$. RQT proves $\tau-t_h\le C_c(1+\log(1/q))$
for sufficiently small $q$. This is a stopping-event theorem and
does not assert that the half-response event occurs first.

Let $g=\nabla n$, $a=Q(1-n)/2$, and
$N=\dot\Theta-ag=-\nabla\mathcal R/4$.
The exact identities, with $\Gamma=\langle g,N\rangle$, are
\begin{equation}\label{an04recv1:face}
 \mathcal R'=-4\|N\|^2-2Q(1-n)\Gamma,
 \qquad
 \Gamma\ge-\frac{2\|N\|^2}{Q(1-n)}
 \quad\hbox{on the upper face}.
\end{equation}
The inequality is the required nonoutward face test, not a proved
bound on the critical path. Since $n\le1/2$ before the stop,
its denominator is strictly positive even if $n$ is not monotone.
The entry is strictly inside the residual cap. At the simultaneous
half-response and residual event the target half response is already
reached; no continuation through that corner is required.

\section{A uniform reduction without a reached angular chart}
Define the vector shape residuals
\[
 e^{(1)}(D)=r(D)-D_1((1-m)e_1-z e_2),\qquad
 e^{(2)}(D)=r(D)-D_2(-h e_1+(1-n)e_2).
\]
Each is orthogonal in its population to the corresponding scalar
$D_p$, component by component. In particular
\begin{equation}\label{an04recv1:shapes}
 \mathcal R=Q_1[(1-m)^2+z^2]+Qh^2+\chi,
 \qquad
 0\le\chi:=\sum_p\E_{P_p}|e^{(p)}|^2
 \le Cq^2(1+S)^2.
\end{equation}
To prove the bound, $f$ is globally $S$-Lipschitz and positively
homogeneous, hence $r$ is $(1+S)$-Lipschitz and $r(0)=0$.
On a cluster with center $\mu e_p$, compare $r(D)$ to
$r(\mu e_p)D_p/\mu$. Their $L^2$ difference is at most
$Cq(1+S)$. Orthogonal projection onto $D_p$ can only improve
this bound. This proof is uniform over all gates and all labels.

Write $g_m=\nabla m$, $g_z=\nabla z$, $g_h=\nabla h$.
For the full feature derivative $J_D$ define
\[
 N_{\rm sh}=\tfrac12\sum_p\E_{P_p}J_D^Te^{(p)}(D).
\]
Exact population gradient flow gives
\begin{equation}\label{an04recv1:four}
 N=\tfrac{Q_1}{2}(1-m)g_m-\tfrac{Q_1}{2}z g_z
             -\tfrac Q2 h g_h+N_{\rm sh}.
\end{equation}
The label-gradient feature bound
$\|J_D^Tv\|^2\le2S|D|^2|v|^2$, followed by
Cauchy--Schwarz in data, gives
\[
 \|N_{\rm sh}\|\le Cq\sqrt S(1+S),\qquad
 \|g\|+\|g_m\|+\|g_z\|+\|g_h\|\le C\sqrt S.
\]
In using Cauchy--Schwarz here one may first use the operator norm
of $J_D$, bounded by $\sqrt{2S}|D|$, so only second data moments
and the $L^2$ shape bound are required.

For a label $U$, put $Q_2=Q$ and define exact Gaussian features
\[
 b_p(U)=\frac{\E_{P_p}D_p(U^TD)_+}{Q_p},\qquad
 v_p(U)=\frac{\E_{P_p}D_pD\,\mathbf1_{U^TD>0}}{Q_p}.
\]
The gradient of output-$i$ regression on cluster $p$ has blocks
$(W_i v_p,e_i b_p)$. Uniformly in $U$,
\begin{equation}\label{an04recv1:features}
 |b_p-(U_p)_+|\le Cq|U|,\qquad
 v_p=\theta_p e_p+O(q),\qquad 0\le\theta_p\le1.
\end{equation}
For the first estimate use the Lipschitz property of the positive
part and $D=\mu e_p+qZ$. For the second choose
$\theta_p=\E_{P_p}D_p^2\mathbf1_{U^TD>0}/Q_p$.
The remaining coordinate is at most
$\sqrt{\E D_p^2\E D_{3-p}^2}/Q_p\le Cq$.
Thus $|v_1^Tv_2|\le Cq$, including every gate boundary.
There is no replacement of a Gaussian gate by a hard center gate.

Define two current-population moments
\begin{equation}\label{an04recv1:moments}
 C_{12}=\int(U_1)_+(U_2)_+\dd l\ge0,\qquad
 T_2=\int W_1W_2|v_2(U)|^2\dd l.
\end{equation}
The latter is signed. The gradient formulas and
\eqref{an04recv1:features} show
\begin{align}
 \langle g,g_m\rangle
  &=\int W_1W_2v_1^Tv_2\dd l=O(qS),\nonumber\\
 \langle g,g_z\rangle
  &=\int(b_1b_2+W_2^2v_1^Tv_2)\dd l=C_{12}+O(qS),\nonumber\\
 \langle g,g_h\rangle&=T_2.\label{an04recv1:grams}
\end{align}
All terms have their physical label masses. In particular
$C_{12}\le S/2$, $|T_2|\le CS$; these bounds do not sign $T_2$.

\begin{theorem}[Two-channel exchange on the stopped critical path]
\label{an04recv1:reduction}
For the above actual trajectory, on $[\tau_0,\tau]$ one has exactly
\begin{equation}\label{an04recv1:gamma}
 \Gamma=-\tfrac12\{Q_1zC_{12}+QhT_2\}+\varepsilon_q,
 \qquad |\varepsilon_q|\le CqS(1+S).
\end{equation}
Here the remainder is explicitly
\[
 \varepsilon_q=\tfrac{Q_1}{2}(1-m)\langle g,g_m\rangle
 -\tfrac{Q_1}{2}z(\langle g,g_z\rangle-C_{12})
 +\langle g,N_{\rm sh}\rangle.
\]
For fixed $c$, as $q\downarrow0$,
\begin{equation}\label{an04recv1:budget}
 \sup_{[\tau_0,\tau]}|\varepsilon_q|
      \le C_c q(1+\log(1/q))^2,
 \qquad
 \int_{\tau_0}^{\tau}|\varepsilon_q|\dd t
      \le C_c q(1+\log(1/q))^3=o(1).
\end{equation}
Also $\sup\chi\le C_cq^2(1+\log(1/q))^2=o(1)$.
\end{theorem}
\begin{proof}
Take the inner product of \eqref{an04recv1:four} with $g$ and
apply \eqref{an04recv1:grams}. Loss decrease bounds $1-m,z,h,1-n$
uniformly by their exact regression inequalities, with $Q,Q_1$
bounded away from zero. This proves \eqref{an04recv1:gamma}.
The global mass identity gives $S(t)\le C(1+t-t_h)$ after the
reached strong half mass. Apply RQT's logarithmic stopping bound
to this identity and integrate. No limiting-system estimate has
been evaluated at a moving time.
\end{proof}

\section{The inequality still missing, and what an exit must pay}
Put $B=Q_1zC_{12}+QhT_2$. Equations
\eqref{an04recv1:face} and \eqref{an04recv1:gamma} give
\begin{equation}\label{an04recv1:missing}
 \boxed{\quad B\le\frac{4\|N\|^2}{Q(1-n)}+2\varepsilon_q\quad}
 \qquad(\mathcal R=r_*^2, n<1/2).
\end{equation}
This is exactly the nonoutward condition. A strict margin or a
local scalar comparison would be needed to turn it into a
no-exit proof. An unsigned $o(1)$ error alone supplies no such margin.
The full dissipation in this test is
\[
 4\|N\|^2=
 \|Q_1(1-m)g_m-Q_1z g_z-Qh g_h+2N_{\rm sh}\|^2.
\]
Its mixed terms cannot be dropped to manufacture lower bounds.

If the residual event occurs before half response, its exact signed
payment is
\begin{equation}\label{an04recv1:exit}
 \int_{\tau_0}^{\tau}Q(1-n)B\dd t
 =r_*^2-\mathcal R(\tau_0)+4\int_{\tau_0}^{\tau}\|N\|^2\dd t
     +2\int_{\tau_0}^{\tau}Q(1-n)\varepsilon_q\dd t.
\end{equation}
In particular its left side is at least
$15r_*^2/16+4\int\|N\|^2-o(1)$.
For sufficiently small $q$ it exceeds $7r_*^2/8$.
At least one of the two \emph{signed integrals}
\[
 \int_{\tau_0}^{\tau}Q(1-n)Q_1zC_{12}\dd t,\qquad
 \int_{\tau_0}^{\tau}Q(1-n)QhT_2\dd t
\]
must therefore exceed $7r_*^2/16$. This conclusion does not
select one of them in advance. It holds only if the residual
event actually occurs. In particular an accumulated fixed-size
exit cannot be attributed solely to Gaussian shape residuals or
to the strong longitudinal mixed term absorbed in $\varepsilon_q$.

The first channel is adverse when the strong-cluster second output
$z$ is positive: shared positive input activation $C_{12}$ couples
the still-unlearned weak target into this cross residual. The second
channel is adverse when the weak-cluster first output $h$ and the
gate-weighted output product $T_2$ have the same sign. It is their
product, not either one's magnitude, that matters. These are
receiver/population correlations in the exact flow, not an
attribution to a fixed ancestral donor cohort.

\paragraph{Two approaches checked.}
First, Cauchy--Schwarz gives
$\Gamma\ge-\|g\|\|N\|$.
For $N\ne0$ this pays \eqref{an04recv1:face} only if
$\|N\|\ge a\|g\|$. Neither the small cap nor loss decrease
gives that lower bound. Completing the square gives only
$\mathcal R'=a^2\|g\|^2-4\|N+ag/2\|^2$, which also
leaves the upper face unsigned.
Second, regression orthogonality removes the shape residuals in
the ordinary data inner product but leaves the gradient Gram terms
in \eqref{an04recv1:grams}. Theorem~\ref{an04recv1:reduction}
pays their approximation errors on the actual stopped interval.
It leaves the two signed products in \eqref{an04recv1:missing}.
Positive semidefiniteness of the Gram matrix supplies no sign or
required size comparison for these mixed entries. The argument
stops here after these two approaches. No actual violation of
\eqref{an04recv1:missing} has been proved.

\paragraph{Scope, constants, and source matching.}
The analytic estimates in \eqref{an04recv1:shapes}--\eqref{an04recv1:grams}
hold directly at finite noise, without chart or ancestry restrictions.
BRW's global mass identity and TBW/RQT's initialized small hit and
stopping bound supply their time interval. Their constants, the
chosen $b,\eta_0$, and the small-noise threshold may depend on fixed
$c,\rho$; $r_*^2=\lambda^2/640000$ does not depend on $c$.
No uniformity as $c\to0$ or $c\to\infty$ is asserted.
FWC's deep-family continuation is not invoked for this critical
entrance. COD's resident estimates remain statements for each fixed
$s$; no COD constant is claimed at $s=t_\eta-t_h$.
The new logarithmic estimates above concern only the physical
regression reduction on the stopped interval, not the scaled probe.
They prove neither half response nor retention of the learned
swing-by advantage. MPB's output-mean mechanism concerns the
lowering of core output; the same evolving cross outputs here can
exchange training residual with weak learning. Their existence
alone does not pay the barrier or the probe-depth inequality.
No numerical trajectory or empirical premise has been used.
All prior increments, foundations, checkpoint and remote files
remain unchanged.
\end{document}
